% From mitthesis package
% Version: 1.01, 2023/06/19
% Documentation: https://ctan.org/pkg/mitthesis
%
% The abstract environment creates all the required headers and footnote. 
% You only need to add the text of the abstract itself.
%
% Approximately 500 words or less; try not to use formulas or special characters
% If you don't want an initial indentation, do \noindent at the start of the abstract

\noindent Neutron stars are compact objects whose hydrostatic equilibrium is determined by the equation of state (EOS) of the dense matter in their interiors, which makes them natural laboratories for nuclear physics at densities inaccessible to terrestrial experiments.
When two neutron stars merge, part of their mass is ejected and, as demonstrated by the gravitational-wave event GW170817 and its associated kilonova, this ejected material is a site of rapid neutron capture (r-process) nucleosynthesis.
The yield of heavy nuclei in the ejecta depends on the properties of neutron-rich nuclei far from stability, including nuclear masses, decay rates, neutron-capture rates, and fission properties, among others, many of which are unknown.
Since nuclear physics plays a crucial role in the evolution of observables, understanding the impact of relying on theoretical nuclear physics results is essential to correctly interpret current and future multimessenger observations.

This work presents a detailed analysis of the impact of different theoretical $\alpha$- and $\beta^{-}$-decay rates on the kilonova evolution.
The nucleosynthesis was calculated using the nuclear reaction network WinNet across an ensemble of initial thermodynamic conditions parametrized by the initial electron fraction ($Y_{e,0}$), the expansion timescale ($\tau$) and the initial entropy per baryon ($s_{0}$).
The kilonova light curves were calculated using current semi-analytical models based on radiative-transfer simulations. 
The results show that the heating rate from radioactive decays deviates from the widely accepted power-law behavior depending on the $\beta^{-}$-decay prescription consistently across the initial conditions studied.
In addition, the set of nuclei responsible for such behavior was identified and the description of their decay chains is presented.

\textcolor{red}{This next paragraph may change as results from GR simulations are being obtained}
The results on the impact of theoretical rates described above are then tested in a merger simulation where the resulting kilonova is now described as a superposition of different ejecta components.
The amount of ejected mass and its velocity are extracted from a general-relativistic simulation which was computed using the Einstein Toolkit framework.



